\begin{proof}[Proof of \cref{obj:thm:all-label-ambiguity}]
\textit{Step 1 (general upper bound).}
For each arm and label, define
\[
\Gamma_{ar}
:=\inf_{c\in[1/(1-\varepsilon),\,1/\varepsilon]}
       \int_{C_r}|1-cp_a(e)|\,H(de),
\qquad
\Gamma_{\mathrm{all}}:=\sum_{r\in\mathcal R_J}
       \bigl(\Gamma_{1r}+\Gamma_{0r}\bigr).
\]
Consider any causal law \(\widetilde P\) in the supremum defining \(D_H(g)\), and set
\[
\widetilde P_{\mathrm{rel}}
:=\mathcal L_{\widetilde P}(R,A,Y).
\]
For a cell with \(q_{ar}>0\), let \(\widetilde F_{ar}\) be its observed outcome law under \(\widetilde P_{\mathrm{rel}}\), and define its contribution to the interval width by
\[
\Delta_{ar}(\widetilde P_{\mathrm{rel}})
:=q_{ar}\int_0^1 Q_{\widetilde F_{ar}}(u)
       \bigl\{Q_{W_{ar}}(u)-Q_{W_{ar}}(1-u)\bigr\}\,du.
\]
Set \(\Delta_{ar}(\widetilde P_{\mathrm{rel}})=0\) when \(q_{ar}=0\). By \cref{obj:lem:fixed-released-law-baseline} and the definitions of the mean endpoints,
\[
\operatorname{len}I(\widetilde P_{\mathrm{rel}};H,g)
=\sum_{r\in\mathcal R_J}
 \bigl\{\Delta_{1r}(\widetilde P_{\mathrm{rel}})
       +\Delta_{0r}(\widetilde P_{\mathrm{rel}})\bigr\}.
\]

Fix a positive-mass cell and \(c\in[1/(1-\varepsilon),1/\varepsilon]\). Since \(0\le Q_{\widetilde F_{ar}}(u)\le1\) almost everywhere, the positive and negative parts give
\[
\begin{aligned}
\int_0^1 Q_{\widetilde F_{ar}}(u)Q_{W_{ar}}(u)\,du
&\le c\int_0^1Q_{\widetilde F_{ar}}(u)\,du
   +\int_0^1\bigl(Q_{W_{ar}}(u)-c\bigr)_+\,du,\\
\int_0^1 Q_{\widetilde F_{ar}}(u)Q_{W_{ar}}(1-u)\,du
&\ge c\int_0^1Q_{\widetilde F_{ar}}(u)\,du
   -\int_0^1\bigl(c-Q_{W_{ar}}(1-u)\bigr)_+\,du.
\end{aligned}
\]
The quantiles at \(u\) and \(1-u\) each have law \(W_{ar}\) under uniform \(u\). Subtraction therefore yields
\[
\frac{\Delta_{ar}(\widetilde P_{\mathrm{rel}})}{q_{ar}}
\le \mathbb E|W_{ar}-c|.
\]
The defining score law of \(E_{ar}\) and \(p_a(e)>0\) give the scaling identity
\[
q_{ar}\mathbb E|W_{ar}-c|
=\int_{C_r}p_a(e)\left|\frac1{p_a(e)}-c\right|H(de)
=\int_{C_r}|1-cp_a(e)|\,H(de).
\]
Taking the infimum shows \(\Delta_{ar}(\widetilde P_{\mathrm{rel}})\le\Gamma_{ar}\). If \(q_{ar}=0\), then \(H(C_r)=0\), since \(p_a(e)\ge\varepsilon>0\) on \(\mathcal E\); hence both sides are zero. Summing over all cells proves
\[
\operatorname{len}I(\widetilde P_{\mathrm{rel}};H,g)
\le\Gamma_{\mathrm{all}},
\qquad D_H(g)\le\Gamma_{\mathrm{all}}.
\]
% lean: CausalSmith.PartialID.UnlinkedPropensityAte.worstCaseAmbiguity_le_sum_cellAbsoluteDeviation

\textit{Step 2 (width of the half-outcome release).}
For the \(P_{\mathrm{rel}}\) in the statement, every positive-mass cell has outcome law Bernoulli\((1/2)\), whose quantile is zero below \(1/2\) and one above \(1/2\), up to a null set. Put \(m_{ar}:=Q_{W_{ar}}(1/2)\) in such a cell. Monotonicity of the quantile and a change of variable give
\[
\begin{aligned}
\frac{\Delta_{ar}(P_{\mathrm{rel}})}{q_{ar}}
&=\int_{1/2}^1Q_{W_{ar}}(u)\,du
  -\int_0^{1/2}Q_{W_{ar}}(u)\,du\\
&=\int_0^1|Q_{W_{ar}}(u)-m_{ar}|\,du
 =\mathbb E|W_{ar}-m_{ar}|.
\end{aligned}
\]
Overlap places \(W_{ar}\), and thus its median \(m_{ar}\), in
\([1/(1-\varepsilon),1/\varepsilon]\). The preceding cellwise bound applies to this fair outcome law for every admissible \(c\), while \(m_{ar}\) attains that bound. The scaling identity consequently gives
\[
\Delta_{ar}(P_{\mathrm{rel}})
=q_{ar}\inf_{c\in[1/(1-\varepsilon),\,1/\varepsilon]}
       \mathbb E|W_{ar}-c|
=\Gamma_{ar}.
\]
The equality also holds in zero-mass cells by the zero-mass argument above. Summing the armwise contributions yields
\[
\operatorname{len}I(P_{\mathrm{rel}};H,g)
=\Gamma_{\mathrm{all}}.
\]
% lean: CausalSmith.PartialID.UnlinkedPropensityAte.halfOutcome_sharpATELength_eq_sum_cellAbsoluteDeviation

\textit{Step 3 (compatible laws and endpoints).}
The cell masses satisfy
\[
\sum_{r\in\mathcal R_J}\sum_{a\in\{0,1\}}q_{ar}
=\int_{\mathcal E}\bigl(p_1(e)+p_0(e)\bigr)\,H(de)=1.
\]
Thus \(P_{\mathrm{rel}}\) is a probability law. For every label \(r\), its defining mixture gives
\[
P_{\mathrm{rel}}(R=r)=q_{1r}+q_{0r}=H(C_r),
\qquad
P_{\mathrm{rel}}(A=1,R=r)=q_{1r}
=\int_{C_r}e\,H(de).
\]
By \cref{obj:prop:compatibility}, there is a compatible full law producing \(P_{\mathrm{rel}}\). Applying \cref{obj:lem:fixed-released-law-baseline} to this released law gives compatible probability laws \(P_-,P_+\) with
\[
\mathbb E_{P_-}[Y_1-Y_0]=\underline\mu_1-\overline\mu_0,
\qquad
\mathbb E_{P_+}[Y_1-Y_0]=\overline\mu_1-\underline\mu_0,
\]
and hence
\[
\mathbb E_{P_+}[Y_1-Y_0]-\mathbb E_{P_-}[Y_1-Y_0]
=\operatorname{len}I(P_{\mathrm{rel}};H,g)
=\Gamma_{\mathrm{all}}.
\]
% lean: CausalSmith.PartialID.UnlinkedPropensityAte.halfOutcome_endpoint_attainment

\textit{Step 4 (supremum and conclusion).}
The set of widths in \cref{obj:def:worst-case-ambiguity} is bounded above by \(\Gamma_{\mathrm{all}}\), by the general upper bound. The compatible law constructed for \(P_{\mathrm{rel}}\) puts its width in that set, and the half-outcome calculation shows that this width equals \(\Gamma_{\mathrm{all}}\). Therefore
\[
D_H(g)=\Gamma_{\mathrm{all}}
=\sum_{r\in\mathcal R_J}\left[
\inf_{c\in[1/(1-\varepsilon),\,1/\varepsilon]}
 \int_{C_r}|1-ce|\,H(de)
+
\inf_{c\in[1/(1-\varepsilon),\,1/\varepsilon]}
 \int_{C_r}|1-c(1-e)|\,H(de)
\right].
\]
Substituting this equality into the attained endpoint gap above proves the asserted gap. The hypotheses also cover the index-set boundary: \(\mathcal E\) contains \(\varepsilon\), so the value \(g(\varepsilon)\) ensures \(\mathcal R_J\) is nonempty.
% lean: CausalSmith.PartialID.UnlinkedPropensityAte.worstCaseAmbiguity_eq
\end{proof}
